\section{Introduction}\label{sec:1}

Two simple sentences make the puzzle visible. \mention{A lot of tourists are in town.} \mention{A lot of money is being spent.} The verb agrees with the \mention{of}-complement (\mention{tourists}, \mention{money}) rather than with the surface head \mention{lot}, which is morphologically singular in both cases. The complement is plural count in the first sentence, singular non-count in the second. The construction lets the relevant feature of the complement reach the verb.

\term{Number-transparent quantificational nouns} is the term used in Huddleston and Pullum's \mention{Cambridge Grammar of the English Language} \autocite*{huddleston-pullum-2002} (henceforth CGEL). Core members include \mention{lot, plenty, lots, bags, heaps, loads, oodles, stacks, remainder, rest, number,} and \mention{couple}, with \mention{majority} and \mention{minority} at the boundary. Membership is small and relatively stable, and members head NPs whose agreement override is, on CGEL's analysis, the default pattern, contrasting with the optional override of ordinary collective nouns. Many other configurations behave similarly: a form mediates between an internal property and a grammatical or interpretive relation, and the property remains accessible through the mediator. The umbrella schema for the family follows.

\begin{quote}
$X$ is \mention{transparent} with respect to property $P$ for relation $R$ when $P$ remains accessible through $X$ for the purposes of $R$.
\end{quote}

The schema is a tool for a particular question: when is calling something \mention{transparent} doing analytical work, and when is it just a label? Sometimes the label genuinely earns its keep; sometimes it doesn't. The \mention{a lot of} construction shows what earning the label looks like.

The schema's substantive content is the \mention{remains accessible through $X$} clause. A case fits only if $X$ has internal structure whose default behaviour would have made $P$ inaccessible to $R$. Cases without that default-conceal background are excluded. Coordination (§\ref{sec:2-3}) is the textbook exclusion: \mention{Bob and Jane are happy} has no head whose morphological number competes for agreement with the conjuncts' features, so there's nothing the construction could have hidden, and calling coordination \mention{transparent} would be a relabel rather than an analysis. Everyday senses of \term{transparent prose} or \term{transparent intentions}, where no specific mediator-property-relation triple is in play, are excluded for a related reason: the schema's slots have no values to bind. The exclusion mechanism is what keeps the schema from cataloguing every use of \mention{transparent} in linguistics, and the rest of the paper is occupied with cases where the slots have non-trivial values.

Once we recognise that \mention{lot, plenty, lots, bags, heaps, loads, oodles, stacks, remainder, rest, number,} and \mention{couple} (in its quantificational sense) form a small, conventionally stable class, we can predict more than the agreement override that motivated the label. Members of the class show characteristic patterns of determiner co-occurrence; they admit partitive \mention{of}-complements in some subgroups but not others; they behave distinctively under the fused-head construction; the count-or-non-count profile of the \mention{of}-complement determines the profile of the whole NP; modification by ordinary adjectives is restricted in characteristic ways. Knowing that a noun belongs to the class lets us predict much of the rest, with the predictions reliable for the core members and graded at the boundary (\mention{majority, minority}, sense-split \mention{couple}). The label is doing work: a small lexical inventory whose membership predicts a cluster of correlated behaviours, attenuated at the edges rather than guaranteed everywhere.

Cases like this we'll call \term{projectibility profiles}. A profile has three components: an \term{observable input} (here, lexical-class membership); a set of \mention{predictions} the input licenses (here, the cluster of agreement, distributional, and modificational behaviours); and a \term{stabilising process} that holds the input together (here, ordinary lexical inheritance). The number-transparent class is a strong example. The predictions are non-trivial; the inventory is stable; the stabilising process is the same one speakers use for any closed grammatical class.

Not every transparency label names a profile this strong. Some name profiles whose predictions are weaker, or whose stabilising process is different, or both. The paper asks which transparency labels do projection-supporting work, and what makes them able to do it.

$X$ (the mediator) can be a lexical item, a partly or fully schematic construction, an intensional context, or a derivationally complex word. $P$ can be a number/count feature, a constituent's meaning, a referent, or a stem's lexical representation. $R$ can be agreement, interpretation, substitution under co-reference, or masked priming. Accessibility itself admits two readings: class-like (core members behave categorically, even where class boundaries are fuzzy) and gradient (compositionality scales, transparency ratings, priming magnitudes).

The argument is constructive. The paper isn't a survey of every use of \term{transparency} in linguistics, nor a paper on extraction, islands, barriers, or phases; extraction-transparency is acknowledged in §\ref{sec:2} but not developed. It isn't a paper on phonological or orthographic transparency, nor on the everyday senses of \mention{transparent explanation} or \mention{transparent terminology}. Referential or \term{de re}/\term{de dicto} transparency receives a controlled mention in §\ref{sec:4} as the boundary case where the schema's mediator-form structure stretches thinnest.

\term{Mediated accessibility} is more specific than \term{clarity} and broader than \term{compositionality}. Clarity is rhetorical: it doesn't single out a mediator or a relation. Compositionality is the limiting case of \term{form-meaning} transparency, where a mediator's parts predict the whole. The umbrella schema is broader: it covers form-meaning cases (constructions, compounds, derived words) and accessibility cases that aren't reducible to whether parts predict wholes. Number transparency is the paradigm of the second type: \mention{lot}'s morphological singular doesn't predict the verb's number, and under the construction it doesn't have to.

Section 2 takes up morphosyntactic transparency, with number-transparent quantificational nouns and transparent free relatives as the paradigm cases. Sections 3 to 5 work through the form-meaning accessibility family: theoretical compositionality \autocite{goldberg-2006-constructions-at-work}, rated compound transparency \autocite{bell-schafer-2016-modelling-transparency}, and time-sensitive lexical processing \autocite{heyer-kornishova-2018-priming-eventually}. Section 6 applies Haspelmath's \autocite*{haspelmath-2010-comparative-concepts} comparative-concept discipline to the schema's cross-linguistic ambitions. Section 7 names the stabilising structures behind each profile and treats Slater's SPC framework \autocite{slater-2015-natural-kindness} and Khalidi's causal-network account \autocite{khalidi-2018-causal-networks} as complementary diagnostics for projection-supporting stability.
